\documentclass[11pt,a4paper]{article}

\usepackage[T1]{fontenc}
\usepackage[utf8]{inputenc}
\usepackage[provide=*,french]{babel}
\usepackage[a4paper,margin=2.3cm]{geometry}
\usepackage{amsmath,amssymb}
\usepackage{booktabs,longtable,array}
\usepackage{hyperref}
\usepackage{enumitem}
\usepackage{listings}
\usepackage{xcolor}

\hypersetup{
  colorlinks=true,
  linkcolor=blue,
  urlcolor=blue,
  pdftitle={Conception technique - Evaluation Potentiel Fongique CHEGD}
}

\lstset{
  basicstyle=\ttfamily\small,
  breaklines=true,
  frame=single,
  backgroundcolor=\color{gray!6}
}

\title{Conception Technique\\
\large Evaluation\_Potentiel\_Fongique\_Interets\_Patrimoniaux\_CHEGD}
\author{Boite Eddy}
\date{\today}

\hbadness=10000
\hfuzz=700pt
\begin{document}
\sloppy
\maketitle

\begin{abstract}
Ce document pr\'esente la conception technique d'un pipeline d'analyse mycologique
destin\'e à l'\'evaluation du potentiel fongique, de l'int\'erêt patrimonial et de la fiabilit\'e de d\'etermination.
La description adopte un cadre d'ing\'enierie reproductible : architecture, contrats de donn\'ees,
m\'ethodes statistiques, robustesse op\'erationnelle et critères de qualit\'e.
L'objectif est de fournir un protocole techniquement auditable, scientifiquement explicitable
et r\'eutilisable pour des campagnes futures.
\end{abstract}
\tableofcontents
\newpage

\section{Structuration IMRaD adapt\'ee}
Dans une perspective de publication scientifique, la lecture du document peut
être organis\'ee selon le sch\'ema IMRaD, adapt\'e à une documentation technique :
\begin{itemize}[leftmargin=1.2cm]
  \item \textbf{Introduction} : positionnement m\'ethodologique, objectifs techniques, p\'erimètre.
  \item \textbf{M\'ethodes} : architecture, modèles de donn\'ees, proc\'edures statistiques, protocoles de validation.
  \item \textbf{R\'esultats} : fichiers g\'en\'er\'es, m\'etriques, tableaux de validation et critères d'assurance qualit\'e (QA).
  \item \textbf{Discussion} : menaces à la validit\'e, limites de transf\'erabilit\'e, recommandations d'\'evolution.
\end{itemize}

Cette grille de lecture vise à renforcer la comparabilit\'e inter-versions et
la r\'eutilisation des livrables dans des contextes acad\'emiques ou institutionnels.

\section{Positionnement m\'ethodologique}
Le pipeline relève d'une approche \textit{data-centric} appliqu\'ee à l'\'ecologie,
où les indicateurs m\'etiers sont construits à partir de transformations successives,
puis consolid\'es par des proc\'edures d'inf\'erence statistique.
Deux objectifs sont poursuivis simultan\'ement :
\begin{itemize}[leftmargin=1.2cm]
  \item \textbf{validit\'e analytique} (qualit\'e des estimations et des tests),
  \item \textbf{validit\'e op\'erationnelle} (stabilit\'e des sorties, exploitabilit\'e d\'ecisionnelle).
\end{itemize}

\section{Cadre th\'eorique des choix statistiques}
Le dispositif repose sur une articulation entre statistique descriptive,
mod\'elisation pr\'edictive et inf\'erence fr\'equentiste. Cette combinaison r\'epond
à une double exigence : d\'ecrire fidèlement les structures observ\'ees et tester
formellement la plausibilit\'e des relations entre variables contextuelles.

Dans ce cadre, la s\'election des m\'etriques privil\'egie :
\begin{itemize}[leftmargin=1.2cm]
  \item la \textbf{comparabilit\'e inter-runs} (mesures stables et interpr\'etables),
  \item la \textbf{sensibilit\'e aux d\'es\'equilibres de classes} (Macro-F1),
  \item la \textbf{capacit\'e explicative} (lecture des effets site/famille/saison),
  \item la \textbf{traçabilit\'e des d\'ecisions} (fichiers de synthèse et journalisation).
\end{itemize}

\section{Objectif technique}
Le document d\'ecrit l'architecture technique du script
\url{scripts/Evaluation_Potentiel_Fongique_Interets_Patrimoniaux_CHEGD.R},
avec un focus sur :
\begin{itemize}[leftmargin=1.2cm]
  \item la chaîne de traitement de donn\'ees,
  \item la logique de calcul des indicateurs m\'etier,
  \item les modèles statistiques utilis\'es,
  \item la robustesse (contrôles, m\'ecanismes de repli, journalisation, QA).
\end{itemize}

\section{Architecture logicielle}
\subsection{Style d'architecture}
Le script suit une architecture \textbf{pipeline batch s\'equentielle}, organis\'ee en modules fonctionnels :
\begin{enumerate}[leftmargin=1.2cm]
  \item Initialisation (packages, configuration embarqu\'ee, journalisation),
  \item Ingestion/normalisation des donn\'ees,
  \item Calculs m\'etier par site,
  \item Module de fiabilit\'e (objectifs statistiques 1 à 4),
  \item G\'en\'eration des exports et visualisations,
  \item Finalisation et journalisation.
\end{enumerate}

\subsection{Principes de conception}
\begin{itemize}[leftmargin=1.2cm]
  \item \textbf{D\'eterminisme} : mêmes entr\'ees + même configuration $\Rightarrow$ mêmes sorties (hors horodatage).
  \item \textbf{Tol\'erance aux variations de donn\'ees} : conversions d\'efensives et diagnostics explicites.
  \item \textbf{Traçabilit\'e dès la conception} : chaque phase significative est journalis\'ee.
  \item \textbf{S\'eparation logique} : ingestion, calcul m\'etier, fiabilit\'e, restitution.
\end{itemize}

\subsection{Points d'entr\'ee}
\begin{itemize}[leftmargin=1.2cm]
  \item \texttt{main()} : orchestration principale.
  \item Bloc \texttt{tryCatch} final : capture d'erreurs globales + log d'\'echec.
\end{itemize}

\subsection{Configuration}
La configuration est int\'egr\'ee via \texttt{get\_embedded\_config()} :
\begin{lstlisting}
protocol_scope = CHEGD pelouses
strict         = TRUE
input_file     = data/donn\'ees_r\'ecoltes_chegd_pelouses.csv
output_dir    = results
output_prefix = EPFIP_CHEGD
columns       = {species, family, date, count, site, reliability}
quality_alert_thresholds = {...}
\end{lstlisting}

Cette section est d\'esormais align\'ee sur les commentaires internes du script,
qui documentent le rôle de chaque paramètre de configuration et la logique de
r\'esolution des chemins relatifs/absolus depuis la racine projet.

\subsection{Mise à jour technique issue des nouveaux commentaires du code}
Les commentaires d\'etaill\'es ajout\'es dans le script renforcent la lisibilit\'e
technique et l'auditabilit\'e. Les points suivants sont explicitement couverts :
\begin{itemize}[leftmargin=1.2cm]
  \item \textbf{Contrats de fonctions} : paramètres, strat\'egie de traitement,
  valeurs de retour et cas limites.
  \item \textbf{R\'esolution d'entr\'ee robuste} : correspondance exacte,
  normalis\'ee, puis rapprochement par distance de Levenshtein.
  \item \textbf{Parsing d\'efensif} : conversions s\'ecuris\'ees des dates et
  num\'eriques, normalisation textuelle et gestion des valeurs manquantes.
  \item \textbf{Mode autonome CSV-only} : lecture exclusive d'un CSV ; les fonctions
  de lecture de classeur externe sont neutralis\'ees (retour \texttt{NULL}).
  \item \textbf{Observabilit\'e} : journalisation structur\'ee de bout en bout avec
  fonctions d\'edi\'ees et traces horodat\'ees.
\end{itemize}

\subsection{Surcharge de configuration}
Bien que la configuration soit embarqu\'ee, le design pr\'evoit une extension simple
vers une surcharge par fichier externe (ex. YAML/JSON) pour :
\begin{itemize}[leftmargin=1.2cm]
  \item les seuils m\'etiers,
  \item les mappings de colonnes,
  \item les chemins d'entr\'ee/sortie,
  \item les paramètres de visualisation.
\end{itemize}

\section{D\'ependances techniques}
\subsection{Packages R}
\begin{itemize}[leftmargin=1.2cm]
  \item \texttt{dplyr}, \texttt{stringr} : transformation/normalisation,
  \item \texttt{ggplot2}, \texttt{gridExtra} : visualisation,
  \item \texttt{MASS} : r\'egression ordinale \texttt{polr},
  \item \texttt{nnet} : r\'egression multinomiale \texttt{multinom}.
\end{itemize}

\subsection{Comportement d'installation}
Les packages requis sont contrôl\'es au d\'emarrage. En cas d'absence, le script s'arrête
avec une erreur explicite et une commande d'installation sugg\'er\'ee (pas d'installation automatique).

\section{Contrats d'entr\'ee et de sortie}
\subsection{Entr\'ees}
Formats support\'es : \texttt{.csv} (s\'eparateur d\'etect\'e automatiquement : \texttt{;}, \texttt{,} ou tabulation).\
Colonnes obligatoires : espèce, famille, date, effectif, site, fiabilit\'e.

\subsection{Sorties}
\begin{itemize}[leftmargin=1.2cm]
  \item R\'epertoire fixe :
  \texttt{results/EPFIP\_CHEGD/}
  \item Organisation automatique en sous-r\'epertoires th\'ematiques :
  \texttt{00\_data\_prepared}, \texttt{10\_indices\_metiers}, \texttt{20\_syntheses},
  \texttt{30\_statistiques\_modeles}, \texttt{40\_qa\_audits},
  \texttt{50\_figures\_metier}, \texttt{60\_figures\_statistiques}
  \item Logs :
  \texttt{logs/EPFIP\_CHEGD\_YYYYMMDD\_HHMM.log}
  \item CSV de synthèse : r\'esum\'es, indicateurs, QA, objectifs fiabilit\'e
  \item Figures PNG/PDF : \texttt{fig1} à \texttt{fig7} + figures objectives.
\end{itemize}

\subsection{Sp\'ecification d\'etaill\'ee des fichiers g\'en\'er\'es}
Le pipeline produit des artefacts en couches : \textit{ingestion},
\textit{indicateurs m\'etier}, \textit{statistiques avanc\'ees}, \textit{visualisations},
\textit{observabilit\'e}. Cette structuration permet d'isoler rapidement une anomalie
dans la chaîne de traitement.

\subsubsection{Couche 1 --- Ingestion et normalisation}
\begin{center}\begin{tabular}{p{0.36\textwidth}p{0.22\textwidth}p{0.34\textwidth}}
\toprule
\textbf{Artefact} & \textbf{Type} & \textbf{Rôle technique} \\
\midrule
\texttt{donnees\_brutes\_nettoyees.csv} & table plate & Point de v\'erit\'e post-nettoyage, base de toutes les agr\'egations. \\
\texttt{resume\_global.csv} & table synthèse & QA global, volum\'etrie, indicateurs de coh\'erence et m\'etadonn\'ees de run. \\
\bottomrule
\end{tabular}\end{center}

\subsubsection{Couche 2 --- Indicateurs m\'etier consolid\'es}
\begin{center}\begin{tabular}{p{0.40\textwidth}p{0.18\textwidth}p{0.34\textwidth}}
\toprule
\textbf{Artefact} & \textbf{Cl\'e primaire logique} & \textbf{Variables critiques attendues} \\
\midrule
potentiel fongique & \texttt{site} & \texttt{potentiel\_score}, \texttt{potentiel\_classe} \\
indice patrimonial & \texttt{site} & \texttt{indice\_patrimonial}, \texttt{classe\_patrimoniale} \\
gradient CHEGD & \texttt{site} & \texttt{gradient\_visite\_*}, \texttt{chegd\_total}, \texttt{chegd\_moyen} \\
indice de repr\'esentativit\'e & \texttt{site} & \texttt{ir\_visite\_*}, \texttt{ir\_moyen} \\
\bottomrule
\end{tabular}\end{center}

\subsubsection{Couche 3 --- Sorties statistiques objectives (fiabilit\'e)}
\begin{center}\begin{tabular}{p{0.23\textwidth}p{0.21\textwidth}p{0.46\textwidth}}
\toprule
\textbf{Pr\'efixe} & \textbf{Objectif} & \textbf{Contenu typique} \\
\midrule
\texttt{stat\_obj1\_*} & O1 & distribution de classes, fr\'equences, proportions, diagnostics descriptifs. \\
\texttt{stat\_obj2\_*} & O2 & scores pond\'er\'es par sch\'ema, $\rho$ de Spearman, overlap top-k. \\
Fichiers modèle O3 & O3 & m\'etriques CV, pr\'edictions, confusion et coefficients. \\
\texttt{stat\_obj4\_*} & O4 & test choisi ($\chi^2$/Fisher), $p$-value, mesure $M=-\log_{10}(p)$. \\
synthèse de s\'election modèle & O3 & synthèse consolid\'ee du meilleur candidat et justification num\'erique. \\
\bottomrule
\end{tabular}\end{center}

\subsubsection{Couche 4 --- Restitution graphique}
\begin{itemize}[leftmargin=1.2cm]
  \item \textbf{Figures de pilotage} : classement sites, gradients, IR, classes patrimoniales.
  \item \textbf{Figures de diagnostic} : distribution fiabilit\'e, dispersion inter-sites, sensibilit\'e aux pond\'erations.
  \item \textbf{Formats} : PNG (consommation rapide), PDF (publication/rapport).
\end{itemize}

\subsubsection{Couche 5 --- Observabilit\'e}
\begin{itemize}[leftmargin=1.2cm]
  \item \texttt{logs/EPFIP\_CHEGD\_*.log} : chronologie d'ex\'ecution, anomalies, dur\'ee, statut final.
  \item \texttt{non\_blocking\_failures.csv} : journal des \'etapes non bloquantes en \'echec (si pr\'esent).
  \item Option recommand\'ee : \textit{manifest} JSON de run (version R, packages, hash input, liste outputs).
\end{itemize}

\subsection{Convention de nommage}
\begin{itemize}[leftmargin=1.2cm]
  \item Pr\'efixe stable m\'etier : \texttt{EPFIP\_CHEGD}.
  \item Noms de fichiers en \texttt{snake\_case} pour faciliter l'automatisation.
  \item Horodatage r\'eserv\'e aux logs et, si besoin, aux exports de diffusion externe.
\end{itemize}

\section{Modèle de donn\'ees technique}
\subsection{Sch\'ema logique minimal des observations}
\begin{center}\begin{tabular}{p{0.22\textwidth}p{0.18\textwidth}p{0.52\textwidth}}
\toprule
\textbf{Champ logique} & \textbf{Type attendu} & \textbf{Rôle dans le pipeline} \\
\midrule
\texttt{species} & caractère & Unit\'e taxonomique de base pour agr\'egations espèce/site. \\
\texttt{family} & caractère & Regroupements taxonomiques et variables explicatives objectif 3. \\
\texttt{date} & date & Saisonnalit\'e, agr\'egations temporelles et qualit\'e de saisie. \\
\texttt{count} & num\'erique & Intensit\'e/abondance utile aux statistiques descriptives et modèles. \\
\texttt{site} & caractère & Cl\'e d'agr\'egation principale des indicateurs. \\
\texttt{reliability} & facteur ordinal & Cible ou poids selon objectifs 1 à 4. \\
\bottomrule
\end{tabular}\end{center}

\subsection{Artefacts interm\'ediaires}
Le pipeline manipule des tables interm\'ediaires standardis\'ees :
\begin{itemize}[leftmargin=1.2cm]
  \item donn\'ees nettoy\'ees (
  \texttt{donnees\_brutes\_nettoyees}),
  \item agr\'egats site/famille/espèce/date,
  \item tables de scores (
  \texttt{potentiel}, \texttt{patrimonial}, \texttt{chegd}),
  \item table consolid\'ee (
  \texttt{combined}) pour export et visualisation.
\end{itemize}

\section{Pipeline technique d\'etaill\'e}
\subsection{Phase A --- Ingestion et normalisation}
\begin{itemize}[leftmargin=1.2cm]
  \item \texttt{resolve\_input\_file} : r\'esolution robuste du chemin (exact, normalis\'e, distance minimale),
  \item \texttt{read\_input\_data} : chargement CSV, d\'etection s\'eparateur CSV,
  \item suppression des colonnes techniques \texttt{...*},
  \item conversion de types :
  \begin{itemize}
    \item \texttt{to\_date\_safe} (dates Excel num\'eriques + formats texte),
    \item \texttt{to\_numeric\_safe} (virgule d\'ecimale).
  \end{itemize}
\end{itemize}

Le pipeline intègre en outre un contrôle de qualit\'e d'entr\'ee export\'e dans
\texttt{qa\_validation\_entree\_*.csv}, ainsi qu'un mode d'ex\'ecution
\texttt{strict}/\texttt{tol\'erant} pilot\'e par la configuration embarqu\'ee.

Les commentaires du code d\'ecrivent \'egalement la d\'etection du s\'eparateur CSV
sur les premières lignes non vides et le nettoyage des colonnes entièrement
vides, afin de fiabiliser les imports issus d'exports h\'et\'erogènes.

\subsubsection{Critères de qualit\'e de donn\'ees}
À l'issue de la phase A, les contrôles suivants sont attendus :
\begin{itemize}[leftmargin=1.2cm]
  \item compl\'etude minimale des variables critiques,
  \item coh\'erence typologique (dates, num\'eriques, facteurs),
  \item traçabilit\'e des exclusions et imputations implicites.
\end{itemize}

\subsection{Phase B --- Calculs m\'etier}
\texttt{build\_site\_level\_metrics} produit 4 tables :
\begin{itemize}[leftmargin=1.2cm]
  \item \texttt{potentiel},
  \item \texttt{patrimonial},
  \item \texttt{chegd},
  \item \texttt{combined} (jointure consolid\'ee).
\end{itemize}

\subsubsection{Potentiel fongique}
Score obtenu par somme pond\'er\'ee de groupes/taxons indicateurs (g\'en\'eriques et sp\'ecifiques).
Classification :
\[
\text{classe}=
\begin{cases}
\text{faible} & \text{si } s\le 10\\
\text{int\'eressant} & \text{si } 10<s<30\\
\text{\'elev\'e} & \text{si } s\ge 30
\end{cases}
\]

\subsubsection{Int\'erêt patrimonial}
Soit $C,H,E,G,D$ les effectifs des 5 groupes CHEGD.
\[
I_{patrimonial}=\max(C,H,E,G,D)
\]
Puis mapping vers classes ordinales (faible/local/r\'egional/national/international).

\subsubsection{Gradient CHEGD}
Comptage d'espèces CHEGD par visite, puis :
\[
CHEGD_{total}=\sum_{v} gradient\_visite_v,
\qquad
CHEGD_{moyen}=\frac{CHEGD_{total}}{N_{visites}}
\]
En mode autonome, les gradients sont calcul\'es directement à partir du CSV et des visites
d\'eriv\'ees des dates observ\'ees. Une exception contrôl\'ee existe pour le jeu canonique
\texttt{donn\'ees\_r\'ecoltes\_chegd\_pelouses.csv}, où un profil de r\'ef\'erence interne peut
recaler certains scores (potentiel/gradient) pour coh\'erence protocolaire.

\subsubsection{Indice de repr\'esentativit\'e (IR)}
Calcul\'e par visite selon la formule interne du script :
\[
IR_{visite}=\max\left(0,1-\frac{gradient\_visite}{nombre\_total}\right)
\]
avec \(nombre\_total = CHEGD_{total}\) au niveau site.
Le pipeline exporte \url{ir_visite_1..5} et \url{ir_moyen}.

\subsection{Phase C --- Module fiabilit\'e (Objectifs 1 à 4)}
Le module est orchestr\'e par \texttt{run\_reliability\_objectives}.

\subsubsection{Contrôles pr\'ealables au module fiabilit\'e}
\begin{itemize}[leftmargin=1.2cm]
  \item V\'erification du nombre minimal d'observations par classe.
  \item V\'erification de la diversit\'e de classes (\'eviter la classe unique).
  \item V\'erification de la variance des variables explicatives principales.
\end{itemize}

\subsubsection{Pr\'eparation des donn\'ees fiabilit\'e}
\begin{itemize}[leftmargin=1.2cm]
  \item normalisation des niveaux : \texttt{Non renseign\'ee}, \texttt{Probable}, \texttt{Certaine},
  \item cr\'eation de variables explicatives : abondance, saison, site, famille,
  \item encodage en facteur ordinal pour la cible.
\end{itemize}

\subsubsection{Objectif 1 --- Descriptif}
Distribution fr\'equence/\% des niveaux de fiabilit\'e.

Sorties : \url{stat_obj1_reliability_distribution.csv} + figure d\'edi\'ee.

\subsubsection{Objectif 2 --- Sch\'emas de pond\'eration}
Trois sch\'emas candidats (S1/S2/S3) pondèrent les niveaux de fiabilit\'e,
puis impactent un score de site.

M\'etriques de comparaison :
\begin{itemize}[leftmargin=1.2cm]
  \item corr\'elation de rang de Spearman $\rho$ entre score de r\'ef\'erence et score pond\'er\'e,
  \item overlap top-5 des sites class\'es.
\end{itemize}
Le meilleur sch\'ema maximise d'abord $\rho$, puis l'overlap.

\paragraph{Sorties techniques attendues (O2)}
\begin{itemize}[leftmargin=1.2cm]
  \item table des scores par site et par sch\'ema,
  \item table comparative des m\'etriques ($\rho$, overlap top-5),
  \item indicateur de robustesse de classement (fort/mod\'er\'e/faible).
\end{itemize}

\subsubsection{Objectif 3 --- S\'election de modèles pr\'edictifs}
\paragraph{Cible}
Variable ordinale \texttt{fiabilite} à 3 niveaux.

\paragraph{Variables explicatives}
\begin{itemize}[leftmargin=1.2cm]
  \item \texttt{abundance},
  \item \texttt{season\_model},
  \item \texttt{site\_model} (top-k + \texttt{Autre}),
  \item \texttt{family\_model} (top-k + \texttt{Autre}).
\end{itemize}

\paragraph{Candidats}
\begin{enumerate}[leftmargin=1.2cm]
  \item \textbf{Ordinal logit} : \texttt{MASS::polr} (hypothèse proportional odds),
  \item \textbf{Multinomial logit} : \texttt{nnet::multinom},
  \item \textbf{Baseline majoritaire} : pr\'ediction de la classe la plus fr\'equente.
\end{enumerate}

\paragraph{Validation}
Validation crois\'ee stratifi\'ee à 5 folds.

\paragraph{M\'etriques}
\begin{itemize}[leftmargin=1.2cm]
  \item Accuracy,
  \item Macro-F1,
  \item Log-loss multiclasses (quand probabilit\'es disponibles).
\end{itemize}

Formules :
\[
\text{Precision}_c=\frac{TP_c}{TP_c+FP_c},
\qquad
\text{Recall}_c=\frac{TP_c}{TP_c+FN_c}
\]
\[
F1_c=\frac{2\,\text{Precision}_c\,\text{Recall}_c}{\text{Precision}_c+\text{Recall}_c},
\qquad
\text{Macro-F1}=\frac{1}{|C|}\sum_{c\in C}F1_c
\]
\[
\text{LogLoss}=-\frac{1}{n}\sum_{i=1}^{n}\log p_{i,y_i}
\]
Le modèle retenu maximise Macro-F1 puis Accuracy.

\paragraph{Sorties techniques attendues (O3)}
\begin{itemize}[leftmargin=1.2cm]
  \item m\'etriques par fold et par modèle,
  \item moyenne/\'ecart-type par m\'etrique,
  \item matrice de confusion agr\'eg\'ee,
  \item fichier de d\'ecision finale de s\'election de modèle.
\end{itemize}

\paragraph{Interpr\'etation scientifique}\leavevmode\\
Le choix d'une m\'etrique prioritaire orient\'ee classes (Macro-F1) limite
l'effet de domination des classes majoritaires et soutient une interpr\'etation
plus \'equilibr\'ee de la performance pr\'edictive dans des jeux potentiellement d\'es\'equilibr\'es.

Cette orientation est coh\'erente avec les standards m\'ethodologiques
en \'ecologie quantitative, où la performance globale doit être examin\'ee
conjointement à la performance par classe pour \'eviter les conclusions biais\'ees.

\subsubsection{\texorpdfstring{Objectif 4 --- Inf\'erence statistique site $\times$ fiabilit\'e}{Objectif 4 --- Inf\'erence statistique site x fiabilit\'e}}
Table de contingence \texttt{site\_label} \texttimes{} \texttt{fiabilite}.

Strat\'egie de test :
\begin{itemize}[leftmargin=1.2cm]
  \item test du $\chi^2$ si conditions d'approximation satisfaites,
  \item sinon test exact de Fisher,
  \item fallback Fisher simul\'e ($B=10000$) si l'exact \'echoue (grandes tables).
\end{itemize}

M\'etrique primaire stock\'ee :
\[
M=-\log_{10}(p)
\]
Plus $M$ est grand, plus l'association est marqu\'ee.

\paragraph{Sorties techniques attendues (O4)}
\begin{itemize}[leftmargin=1.2cm]
  \item table de contingence \texttt{site\_label} $\times$ \texttt{fiabilite},
  \item type de test effectivement utilis\'e,
  \item statistique de test, $p$-value et mesure $M$,
  \item drapeau de fallback (exact/simul\'e) pour traçabilit\'e m\'ethodologique.
\end{itemize}

\subsection{Phase D --- Visualisation}
Fonctions d\'edi\'ees :
\begin{itemize}[leftmargin=1.2cm]
  \item \texttt{build\_site\_dashboard} (3 panneaux),
  \item \texttt{build\_scatter\_positionnement},
  \item \texttt{build\_potentiel\_breakdown},
  \item \texttt{build\_chegd\_by\_visit},
  \item \texttt{build\_classes\_heatmap},
  \item \texttt{build\_reliability\_levels\_plot},
  \item \texttt{build\_ir\_by\_visit\_plot},
  \item figures sp\'ecifiques objectifs 1 à 4.
\end{itemize}

\subsection{Contrats graphiques}
Les graphiques doivent respecter les principes suivants :
\begin{itemize}[leftmargin=1.2cm]
  \item lisibilit\'e sur fond clair, palettes contrast\'ees, l\'egendes explicites,
  \item titres/axes homogènes en français,
  \item reproductibilit\'e des dimensions d'export (PNG/PDF).
\end{itemize}

\section{Mode autonome (sans d\'ependance classeur externe)}
Le script opère en mode autonome : les fonctions de lecture/r\'esolution d'un
classeur de r\'ef\'erence externe sont neutralis\'ees et retournent \texttt{NULL}.

Cons\'equences techniques :
\begin{itemize}[leftmargin=1.2cm]
  \item les m\'etriques site sont calcul\'ees directement à partir du CSV d'entr\'ee ;
  \item aucun alignement automatique sur un classeur externe n'est appliqu\'e ;
  \item les sorties restent entièrement d\'eterministes pour un même input et une même configuration.
\end{itemize}

\section{Ex\'ecution stricte/tol\'erante et \'etapes non bloquantes}
Le paramètre \texttt{strict} de la configuration embarqu\'ee contrôle le comportement en cas d'anomalie QA :
\begin{itemize}[leftmargin=1.2cm]
  \item \texttt{strict = TRUE} : arrêt du pipeline sur contrôles bloquants.
  \item \texttt{strict = FALSE} : warning explicite et poursuite contrôl\'ee.
\end{itemize}

La g\'en\'eration des figures et le module fiabilit\'e sont encapsul\'es via
\texttt{run\_non\_blocking\_step()}. En cas d'\'echec local, le pipeline continue son exécution, journalise l'erreur et exporte \texttt{non\_blocking\_failures.csv}.

\section{Logging et observabilit\'e}
\subsection{Fonctions d\'edi\'ees}
\begin{itemize}[leftmargin=1.2cm]
  \item \texttt{setup\_logging}, \texttt{log\_info}, \texttt{log\_warning\_msg}, \texttt{log\_error},
  \item \texttt{log\_section}, \texttt{log\_header}, \texttt{log\_data\_summary}, \texttt{log\_footer}.
\end{itemize}

\subsection{Comportement}
\begin{itemize}[leftmargin=1.2cm]
  \item duplication console + fichier,
  \item messages en français,
  \item \'etapes majeures horodat\'ees,
  \item en-tête + r\'esum\'e de donn\'ees + pied de page (dur\'ee/chemin).
\end{itemize}

\subsection{Év\'enements cl\'es attendus dans le log}
\begin{itemize}[leftmargin=1.2cm]
  \item d\'emarrage et chargement de configuration,
  \item \'etat des donn\'ees en entr\'ee (volum\'etries, colonnes d\'etect\'ees),
  \item passage de chaque phase A/B/C/D,
  \item sorties produites (chemins de fichiers),
  \item erreurs bloquantes et stack simplifi\'ee,
  \item dur\'ee totale d'ex\'ecution.
\end{itemize}

\section{Robustesse et gestion d'erreurs}
\subsection{Niveaux de protection}
\begin{itemize}[leftmargin=1.2cm]
  \item validation de sch\'ema d'entr\'ee,
  \item conversion d\'efensive des types,
  \item \texttt{tryCatch} local sur modèles/tests sensibles,
  \item \texttt{tryCatch} global autour de \texttt{main()}.
\end{itemize}

\subsection{Fallbacks notables}
\begin{itemize}[leftmargin=1.2cm]
  \item fallback Fisher simul\'e en inf\'erence,
  \item baseline majoritaire si modèles indisponibles,
  \item gestion classe unique / donn\'ees insuffisantes pour l'objectif 3.
\end{itemize}

\subsection{Politique de s\'ev\'erit\'e des anomalies}
\begin{center}\begin{tabular}{>{\raggedright\arraybackslash}p{0.2\textwidth}>{\raggedright\arraybackslash}p{0.32\textwidth}>{\raggedright\arraybackslash}p{0.4\textwidth}}
\toprule
\textbf{Niveau} & \textbf{Exemple} & \textbf{Action système} \\
\midrule
INFO & Volum\'etrie et progression & Trace sans interruption \\
ALERTE & Valeurs partiellement manquantes/non standard & Poursuite avec normalisation ou exclusion cibl\'ee \\
ERREUR & Colonnes obligatoires absentes, lecture impossible & Arrêt propre avec message explicite \\
\bottomrule
\end{tabular}\end{center}

\section{Contrôles qualit\'e et non-r\'egression}
\subsection{Indicateurs d'assurance qualit\'e (QA) export\'es}
Dans \texttt{resume\_global.csv} :
\begin{itemize}[leftmargin=1.2cm]
  \item \texttt{coherence\_chegd\_gradients\_ok},
  \item \texttt{nb\_sites\_incoherents\_chegd},
  \item \texttt{ir\_moyen\_global},
  \item m\'etriques volum\'etriques (nb lignes, espèces, familles, sites, dates).
\end{itemize}

\subsection{Critères techniques d'acceptation}
\begin{enumerate}[leftmargin=1.2cm]
  \item ex\'ecution sans erreur fatale,
  \item g\'en\'eration complète des sorties CSV/figures,
  \item log horodat\'e pr\'esent dans \texttt{logs/},
  \item coh\'erence gradients CHEGD valid\'ee.
\end{enumerate}

\subsection{Strat\'egie de tests recommand\'ee}
\begin{itemize}[leftmargin=1.2cm]
  \item \textbf{Tests unitaires} : fonctions de conversion et formules m\'etier (potentiel, CHEGD, IR).
  \item \textbf{Tests d'int\'egration} : ex\'ecution de bout en bout sur un jeu r\'eduit connu.
  \item \textbf{Tests de non-r\'egression} : comparaison d'exports cl\'es avec des snapshots de r\'ef\'erence.
  \item \textbf{Tests de robustesse} : jeux volontairement bruit\'es (dates ambiguës, valeurs manquantes, classes rares).
\end{itemize}

\subsection{Matrice de validation des artefacts}
\begin{center}\begin{tabular}{p{0.21\textwidth}p{0.34\textwidth}p{0.37\textwidth}}
\toprule
\textbf{Bloc} & \textbf{Règle de validation} & \textbf{Preuve attendue} \\
\midrule
Ingestion & Colonnes critiques pr\'esentes et typ\'ees & \texttt{donnees\_brutes\_nettoyees.csv} + warnings cibl\'es dans log \\
Indicateurs m\'etier & 1 ligne par site, pas de duplication cl\'e & tables \texttt{*\_par\_site.csv} coh\'erentes \\
CHEGD/IR & somme gradients visite coh\'erente avec total & QA indicateur de coh\'erence CHEGD \\
Objectif 2 & classement stable sous pond\'erations & \texttt{stat\_obj2\_*.csv} avec $\rho$ et overlap \\
Objectif 3 & modèle gagnant justifi\'e num\'eriquement & synthèse de s\'election modèle \\
Objectif 4 & type de test trac\'e et interpr\'etable & \texttt{stat\_obj4\_*.csv} + log test \\
Observabilit\'e & run horodat\'e complet & fichier \texttt{logs/EPFIP\_CHEGD\_*.log} \\
\bottomrule
\end{tabular}\end{center}

\subsection{Critères de stabilit\'e statistique recommand\'es}
Pour une exploitation en contexte institutionnel, il est recommand\'e de viser :
\begin{itemize}[leftmargin=1.2cm]
  \item une variabilit\'e inter-fold mod\'er\'ee des m\'etriques O3 (\'ecart-type limit\'e),
  \item une corr\'elation de rang O2 \'elev\'ee (id\'ealement $\rho \geq 0.85$),
  \item une coh\'erence des conclusions O4 entre test principal et fallback simul\'e,
  \item une documentation explicite de toute d\'eviation m\'ethodologique dans le log.
\end{itemize}

\subsection{Menaces à la validit\'e}
\begin{itemize}[leftmargin=1.2cm]
  \item \textbf{Validit\'e interne} : d\'epend de la qualit\'e des transformations de variables et de la stabilit\'e des mappings.
  \item \textbf{Validit\'e externe} : transf\'erabilit\'e limit\'ee si nouveaux territoires ou protocoles d'\'echantillonnage très diff\'erents.
  \item \textbf{Validit\'e de conclusion} : prudence requise en cas d'effectifs faibles par classe pour l'inf\'erence site $\times$ fiabilit\'e.
\end{itemize}

\section{Complexit\'e et performances (qualitatives)}
\begin{itemize}[leftmargin=1.2cm]
  \item Agr\'egations principales : complexit\'e proche de $O(n)$ à $O(n\log n)$ selon tris/groupements,
  \item Objectif 3 : coût dominant, approximativement proportionnel à
  \(k \times m\) (folds \texttimes{} modèles) multipli\'e par le coût d'entraînement,
  \item Visualisations : coût lin\'eaire en nombre de points/agr\'egats.
\end{itemize}

\subsection{Pistes d'optimisation}
\begin{itemize}[leftmargin=1.2cm]
  \item limiter les copies de tables interm\'ediaires volumineuses,
  \item mutualiser les agr\'egations redondantes,
  \item m\'emoriser certains objets d\'eriv\'es (facteurs top-k) entre objectifs,
  \item parall\'eliser la validation crois\'ee si n\'ecessaire (selon environnement d'ex\'ecution).
\end{itemize}

\section{S\'ecurit\'e, conformit\'e et exploitation}
\subsection{S\'ecurit\'e des donn\'ees}
\begin{itemize}[leftmargin=1.2cm]
  \item pas de secret applicatif requis pour l'ex\'ecution standard,
  \item privil\'egier des chemins locaux contrôl\'es pour les donn\'ees sources,
  \item limiter la diffusion des logs si ceux-ci contiennent des m\'etadonn\'ees sensibles.
\end{itemize}

\subsection{Conformit\'e et bonnes pratiques}
\begin{itemize}[leftmargin=1.2cm]
  \item conserver la provenance des donn\'ees (source, date d'extraction),
  \item documenter les transformations non triviales,
  \item archiver le triplet \textit{input + outputs + log} pour auditabilit\'e.
\end{itemize}

\subsection{Proc\'edure op\'eratoire standard (POS)}
\begin{enumerate}[leftmargin=1.2cm]
  \item V\'erifier pr\'esence du fichier d'entr\'ee et structure minimale des colonnes.
  \item Ex\'ecuter le script dans l'environnement R cible.
  \item Contrôler le log et les indicateurs QA de \texttt{resume\_global.csv}.
  \item Diffuser les exports valid\'es vers l'\'equipe m\'etier.
  \item Archiver les artefacts de run avec identifiant de campagne.
\end{enumerate}

\section{R\'eplicabilit\'e et transparence scientifique}
Afin de favoriser la r\'eplication des r\'esultats, chaque ex\'ecution devrait documenter :
\begin{itemize}[leftmargin=1.2cm]
  \item la version de R et des packages,
  \item le hachage ou l'empreinte du jeu d'entr\'ee,
  \item les paramètres de configuration effectifs,
  \item les sorties consolid\'ees (CSV, figures, log) associ\'ees au même identifiant de run.
\end{itemize}

\subsection{Sp\'ecification d'un manifeste de run (recommand\'ee)}
Format propos\'e : \texttt{run\_manifest.json} stock\'e à la racine des r\'esultats.
Champs minimaux sugg\'er\'es :
\begin{itemize}[leftmargin=1.2cm]
  \item \texttt{run\_id}, \texttt{timestamp\_start}, \texttt{timestamp\_end},
  \item \texttt{r\_version}, \texttt{packages} (versions),
  \item \texttt{input\_file}, \texttt{input\_hash},
  \item \texttt{config\_effective},
  \item \texttt{outputs\_generated} (liste exhaustive),
  \item \texttt{qa\_summary} (coh\'erence CHEGD, nombre d'alertes, statut final).
\end{itemize}

Ce manifeste facilite les audits comparatifs inter-campagnes et la reproductibilit\'e
scientifique au-delà du seul environnement local d'ex\'ecution.

\section{Évolutions techniques recommand\'ees}
\begin{itemize}[leftmargin=1.2cm]
  \item externaliser les seuils et poids m\'etier dans un bloc param\'etrable,
  \item factoriser les fonctions communes entre script principal et alias historique,
  \item ajouter une suite de tests unitaires ciblant les formules m\'etier,
  \item versionner les modèles/statistiques de fiabilit\'e avec m\'etadonn\'ees de run.
\end{itemize}

\section{Protocole de base}
Ce pipeline impl\'emente le protocole d'\'evaluation d\'ecrit par Yann Sellier \& co :

\begin{quote}
  \textit{Sellier, Y. et coll.} (s.d.). Évaluation du potentiel fongique et de l'int\'erêt patrimonial des pelouses.
  Ressource en ligne : \url{https://www.mycofrance.fr/wp-content/uploads/2025/05/Bull.-SMF-131-1-2-Y-Sellier-et-coll.pdf}.
\end{quote}

Le protocole original sp\'ecifie : 
\begin{itemize}[leftmargin=1.2cm]
  \item Les cinq groupes CHEGD : Clavaria, Hygrocybe, Entoloma, Geoglossaceae, Dermoloma.
  \item Les pond\'erations de potentiel fongique par espèce/groupe indicateur.
  \item L'indice patrimonial comme maximum des effectifs CHEGD.
  \item L'indice de repr\'esentativit\'e relatif au gradient CHEGD par visite.
\end{itemize}

\section{R\'ef\'erences bibliographiques (style harmonis\'e)}
\begin{itemize}[leftmargin=1.2cm]
  \item Hastie, T., Tibshirani, R., \& Friedman, J. (2009). \textit{The elements of statistical learning} (2nd ed.). Springer.
  \item ISO. (s.d.). \textit{ISO 8000 series -- Data quality}. International Organization for Standardization.
  \item James, G., Witten, D., Hastie, T., \& Tibshirani, R. (2021). \textit{An introduction to statistical learning} (2nd ed.). Springer.
  \item Sellier, Y., et coll. (s.d.). \textit{Évaluation du potentiel fongique et de l'int\'erêt patrimonial des pelouses}. Mycofrance. \url{https://www.mycofrance.fr/wp-content/uploads/2025/05/Bull.-SMF-131-1-2-Y-Sellier-et-coll.pdf}.
  \item Wickham, H., \& Grolemund, G. (2017). \textit{R for data science}. O'Reilly.
\end{itemize}
% ============================================================================================================================
\section*{Contact et versionnement}
\addcontentsline{toc}{section}{Contact et versionnement}
% ============================================================================================================================

\begin{table}[htbp]
  \centering
  \begin{tabular}{ll}
    \toprule
    \textbf{Auteur}          & Eddy Boite \\
    \textbf{D\'epôt GitHub}    & \texttt{eddyboite-nat/myco\_apps\_releases} \\
    \textbf{Branche}         & \texttt{main} \\
    \textbf{Version script}  & 1.1 \\
    \textbf{Date de ce document} & 9 juillet 2026 \\
    \bottomrule
  \end{tabular}
\end{table}

Pour toute \'evolution du script, ouvrir une issue dans le d\'epôt GitHub
ou documenter les changements dans le journal \texttt{CHANGELOG.md} du projet.

% ============================================================================================================================
\end{document}
